% ===========================================================================
\section{结构与满足}\label{sec:semantics}

语法规定符号串的合法构造，语义则规定符号的解释。一阶语言中的函数符号、
谓词符号与常元符号本身只是名称，其含义由结构指定。

\begin{definition}[结构]
一阶语言的一个\emph{结构} $\M=(\Omega;\ \ldots)$ 由以下部分组成：
\begin{itemize}[leftmargin=1.3em,itemsep=1pt]
  \item \keyword{论域} $\Omega$：一个\emph{非空}集合（universe）；
  \item 对每个 $k$ 元函数符号 $f_i^k$，指定函数 $f_i^k:\Omega^k\to\Omega$；
  \item 对每个 $k$ 元谓词符号 $P_i^k$，指定关系 $P_i^k\subseteq\Omega^k$；
  \item 对每个常元符号 $c_i$，指定元素 $c_i\in\Omega$。
\end{itemize}
\end{definition}

\begin{remark}[论域非空的理由]
若允许 $\Omega=\varnothing$，则 $\forall x\,(x=x)$ 在空论域上空虚地为真，
而 $\exists x\,(x=x)$ 为假。相应的推理链
\[
  \forall x\,(x=x)\ \longrightarrow\ x=x\ \longrightarrow\ \exists x\,(x=x)
\]
在空论域上失效：第一步是全称消去，第二步是存在引入，而后者要求论域中
确实存在可作见证的元素。即由全称命题推不出相应的存在命题。
允许空论域的逻辑体系称为\emph{自由逻辑}（free logic）；本课程一律假定论域非空。
\end{remark}

\block{解释的递归定义}
给定结构 $\M$，项的解释与公式的真值 $\sem{\cdot}_{\M}$ 递归定义如下：
\begin{align*}
  \sem{c_i}_{\M} &= c_i,\\
  \sem{f_i^k(t_1,\ldots,t_k)}_{\M} &= f_i^k\bigl(\sem{t_1}_{\M},\ldots,\sem{t_k}_{\M}\bigr),\\
  \sem{P_i^k(t_1,\ldots,t_k)}_{\M} &= 1
    \ \Longleftrightarrow\ \bigl(\sem{t_1}_{\M},\ldots,\sem{t_k}_{\M}\bigr)\in P_i^k,\\
  \sem{\forall x\,\psi}_{\M} &= \min_{v\in\Omega}\ \sem{\psi[v/x]}_{\M}.
\end{align*}
末行右端的 $v/x$ 需要说明：$v$ 是论域中的元素，并非语言中的符号，直接代入
不合语法。处理方式是\keyword{扩张语言}：为每个 $v\in\Omega$ 引入一个以它为解释的
常元符号，从而 $\psi[v/x]$ 仍是合法代入。等价的常见处理是使用变量赋值
$s[x\mapsto v]$。

\begin{remark}[等号的解释]
等号不作为普通的二元谓词符号处理，其解释固定为论域上的相等关系：
\[
  \sem{t_1=t_2}_{\M}=1 \ \Longleftrightarrow\ \sem{t_1}_{\M}=\sem{t_2}_{\M}.
\]
\end{remark}

\begin{definition}[满足、模型、逻辑后承]\label{def:satisfaction}
\leavevmode
\begin{itemize}[leftmargin=1.3em,itemsep=2pt]
  \item 对\emph{句子} $\varphi$：$\M\models\varphi$ 当且仅当 $\sem{\varphi}_{\M}=1$；
        此时称 $\M$ 为 $\varphi$ 的一个\keyword{模型}（model）。
  \item 对一般\emph{公式} $\varphi$（自由变元为 $x_1,\ldots,x_n$）：
        $\M\models\varphi$ 定义为 $\M\models\forall x_1\forall x_2\cdots\forall x_n\,\varphi$，
        即取\emph{全称闭包}。按此定义，$\M\models\varphi$ 同样不依赖变量赋值。
  \item 对公式集 $\Gamma$：$\M\models\Gamma$ 表示 $\M\models\varphi$ 对所有 $\varphi\in\Gamma$ 成立。
  \item $\models\varphi$（$\varphi$ \emph{有效}）表示 $\M\models\varphi$ 对所有结构 $\M$ 成立。
  \item $T\models\varphi$（\emph{逻辑后承}）表示：对所有结构 $\M$，$\M\models T$ 蕴含 $\M\models\varphi$。
\end{itemize}
\end{definition}

句子与一般公式分开定义的原因：句子的真值本就不依赖变量赋值，而含自由变元的
公式须先以全称量词约束全部自由变元，满足关系才有确定含义。
